% GENERATED FILE. Edit canonical lessons, then run npm run generate:books.
% canonical-source-hash: fnv1a64:bf64b7022fd3285b
% canonical-lessons: HI-C364-visesajna, HI-C364-vidvan, HI-C364-lekh, HI-C364-nibandh, HI-C364-chatna

\chapter{Ask an Expert}
\label{ch:hi-ask-an-expert}
\hlchaptermodality{364}

\begin{chapteropening}
\textbf{By the end of this chapter:} \emph{I can say who knows best and where they said it.}

Complete the last lesson of chapter 364: I can say who knows best and where they said it.
\end{chapteropening}

\section[\texorpdfstring{\hi{विशेषज्ञ} (viśeṣajña)}{विशेषज्ञ (viśeṣajña)}]{\hi{विशेषज्ञ} (viśeṣajña) — an expert}

\label{lesson:HI-C364-visesajna}



\begin{quote}
Before the new one: say the Hindi for to compare, then the Hindi for to assess.
\end{quote}

\subsection*{You'll want to know: \hi{विशेषज्ञ}}
\textbf{\hi{विशेषज्ञ}} — \emph{viśeṣajña} — \textquotedblleft{}an expert\textquotedblright{}.

\begin{quote}
\hi{हमें} \hi{एक} \hi{विशेषज्ञ} \hi{से} \hi{मिलना} \hi{चाहिए।}
\end{quote}

\begin{quote}
\emph{hamẽ ek viśeṣajña se milnā cāhiye.}
\end{quote}

\begin{quote}
We should see an expert.
\end{quote}

Masculine or feminine, by the person meant.

\begin{cousinweb}
\textbf{\hi{विशेषज्ञ}} (\emph{viśeṣajña}) is a tatsama, borrowed from Sanskrit \emph{viśeṣajña}, \textquotedblleft{}knowing the differences\textquotedblright{} (Monier-Williams).
\end{cousinweb}

\begin{grammarlens}[title={What you've built}]
One more word for giving reasons and opinions.
\end{grammarlens}

\subsection*{Your turn}
\begin{itemize}
\raggedright
  \item \emph{Say it:} \emph{viśeṣajña}
  \item \emph{Say it:} \emph{viśeṣajña}, once more
  \item \emph{Say it:} \emph{hamẽ ek viśeṣajña se milnā cāhiye.} — We should see an expert.
  \item \emph{Recall:} say the Hindi for to compare, then the Hindi for to assess, then say \emph{viśeṣajña} again
\end{itemize}

\begin{tcolorbox}[breakable,colback=teal!4,colframe=teal!35!black,title={Before you move on}]
What does \hi{विशेषज्ञ} mean? (\textquotedblleft{}An expert\textquotedblright{}.) Say it once more.
\end{tcolorbox}
\section[\texorpdfstring{\hi{विद्वान} (vidvān)}{विद्वान (vidvān)}]{\hi{विद्वान} (vidvān) — a scholar, a learned person}

\label{lesson:HI-C364-vidvan}



\begin{quote}
Before the new one: say the Hindi for to assess, then the Hindi for an expert.
\end{quote}

\subsection*{You'll want to know: \hi{विद्वान}}
\textbf{\hi{विद्वान}} — \emph{vidvān} — \textquotedblleft{}a scholar, a learned person\textquotedblright{}.

\begin{quote}
\hi{मेरे} \hi{दादा} \hi{इतिहास} \hi{के} \hi{बड़े} \hi{विद्वान} \hi{थे।}
\end{quote}

\begin{quote}
\emph{mere dādā itihās ke baṛe vidvān the.}
\end{quote}

\begin{quote}
My grandfather was a great scholar of history.
\end{quote}

Masculine or feminine, by the person meant.

\begin{cousinweb}
\textbf{\hi{विद्वान}} (\emph{vidvān}) is a tatsama, borrowed from Sanskrit \emph{vidvas}, \textquotedblleft{}knowing, wise\textquotedblright{} (Monier-Williams).
\end{cousinweb}

\begin{grammarlens}[title={What you've built}]
One more word for giving reasons and opinions.
\end{grammarlens}

\subsection*{Your turn}
\begin{itemize}
\raggedright
  \item \emph{Say it:} \emph{vidvān}
  \item \emph{Say it:} \emph{vidvān}, once more
  \item \emph{Say it:} \emph{mere dādā itihās ke baṛe vidvān the.} — My grandfather was a great scholar of history.
  \item \emph{Recall:} say the Hindi for to assess, then the Hindi for an expert, then say \emph{vidvān} again
\end{itemize}

\begin{tcolorbox}[breakable,colback=teal!4,colframe=teal!35!black,title={Before you move on}]
What does \hi{विद्वान} mean? (\textquotedblleft{}A scholar, a learned person\textquotedblright{}.) Say it once more.
\end{tcolorbox}
\section[\texorpdfstring{\hi{लेख} (lekh)}{लेख (lekh)}]{\hi{लेख} (lekh) — an article (in a paper)}

\label{lesson:HI-C364-lekh}



\begin{quote}
Before the new one: say the Hindi for an expert, then the Hindi for a scholar.
\end{quote}

\subsection*{You'll want to know: \hi{लेख}}
\textbf{\hi{लेख}} — \emph{lekh} — \textquotedblleft{}an article (in a paper)\textquotedblright{}.

\begin{quote}
\hi{मैंने} \hi{अख़बार} \hi{में} \hi{उसका} \hi{लेख} \hi{पढ़ा।}
\end{quote}

\begin{quote}
\emph{maiṁne akhbār meṁ uskā lekh paṛhā.}
\end{quote}

\begin{quote}
I read his article in the newspaper.
\end{quote}

Masculine.

\begin{cousinweb}
\textbf{\hi{लेख}} (\emph{lekh}) is a tatsama, borrowed from Sanskrit \emph{lekha}, \textquotedblleft{}a writing, a document\textquotedblright{} (Monier-Williams).
\end{cousinweb}

\begin{grammarlens}[title={What you've built}]
One more word for giving reasons and opinions.
\end{grammarlens}

\subsection*{Your turn}
\begin{itemize}
\raggedright
  \item \emph{Say it:} \emph{lekh}
  \item \emph{Say it:} \emph{lekh}, once more
  \item \emph{Say it:} \emph{maiṁne akhbār meṁ uskā lekh paṛhā.} — I read his article in the newspaper.
  \item \emph{Recall:} say the Hindi for an expert, then the Hindi for a scholar, then say \emph{lekh} again
\end{itemize}

\begin{tcolorbox}[breakable,colback=teal!4,colframe=teal!35!black,title={Before you move on}]
What does \hi{लेख} mean? (\textquotedblleft{}An article (in a paper)\textquotedblright{}.) Say it once more.
\end{tcolorbox}
\section[\texorpdfstring{\hi{निबंध} (nibandh)}{निबंध (nibandh)}]{\hi{निबंध} (nibandh) — an essay}

\label{lesson:HI-C364-nibandh}



\begin{quote}
Before the new one: say the Hindi for a scholar, then the Hindi for an article.
\end{quote}

\subsection*{You'll want to know: \hi{निबंध}}
\textbf{\hi{निबंध}} — \emph{nibandh} — \textquotedblleft{}an essay\textquotedblright{}.

\begin{quote}
\hi{मीरा} \hi{ने} \hi{विज्ञान} \hi{पर} \hi{एक} \hi{निबंध} \hi{लिखा।}
\end{quote}

\begin{quote}
\emph{Mīrā ne vijñān par ek nibandh likhā.}
\end{quote}

\begin{quote}
Mira wrote an essay on science.
\end{quote}

Masculine.

\begin{cousinweb}
\textbf{\hi{निबंध}} (\emph{nibandh}) is a tatsama, borrowed from Sanskrit \emph{nibandha}, \textquotedblleft{}binding; a literary composition\textquotedblright{} (Monier-Williams).
\end{cousinweb}

\begin{grammarlens}[title={What you've built}]
One more word for giving reasons and opinions.
\end{grammarlens}

\subsection*{Your turn}
\begin{itemize}
\raggedright
  \item \emph{Say it:} \emph{nibandh}
  \item \emph{Say it:} \emph{nibandh}, once more
  \item \emph{Say it:} \emph{Mīrā ne vijñān par ek nibandh likhā.} — Mira wrote an essay on science.
  \item \emph{Recall:} say the Hindi for a scholar, then the Hindi for an article, then say \emph{nibandh} again
\end{itemize}

\begin{tcolorbox}[breakable,colback=teal!4,colframe=teal!35!black,title={Before you move on}]
What does \hi{निबंध} mean? (\textquotedblleft{}An essay\textquotedblright{}.) Say it once more.
\end{tcolorbox}
\section[\texorpdfstring{\hi{छाँटना} (chā̃ṭnā)}{छाँटना (chā̃ṭnā)}]{\hi{छाँटना} (chā̃ṭnā) — to sort out, to pick out}

\label{lesson:HI-C364-chatna}



\begin{quote}
Before the new one: say the Hindi for an article, then the Hindi for an essay.
\end{quote}

\subsection*{You'll want to know: \hi{छाँटना}}
\textbf{\hi{छाँटना}} — \emph{chā̃ṭnā} — \textquotedblleft{}to sort out, to pick out\textquotedblright{}.

\begin{quote}
\hi{माँ} \hi{अच्छे} \hi{आम} \hi{छाँट} \hi{रही} \hi{है।}
\end{quote}

\begin{quote}
\emph{māṁ acche ām chā̃ṭ rahī hai.}
\end{quote}

\begin{quote}
Mother is picking out the good mangoes.
\end{quote}

Past, with \hi{ने} (he or she …): \hi{उसने} \hi{अच्छे} \hi{फल} \hi{छाँटे}, \emph{usne acche phal chā̃ṭe}.

\begin{cousinweb}
\textbf{\hi{छाँटना}} (\emph{chā̃ṭnā}) is inherited from \emph{*chaṇṭ}, \textquotedblleft{}scatter\textquotedblright{}, a form CDIAL reconstructs from its descendants (CDIAL 4970). Hindi uses it for sifting and picking out.
\end{cousinweb}

\begin{grammarlens}[title={What you've built}]
One more word for giving reasons and opinions.
\end{grammarlens}

\subsection*{Your turn}
\begin{itemize}
\raggedright
  \item \emph{Say it:} \emph{chā̃ṭnā}
  \item \emph{Say it:} \emph{chā̃ṭnā}, once more
  \item \emph{Say it:} \emph{māṁ acche ām chā̃ṭ rahī hai.} — Mother is picking out the good mangoes.
  \item \emph{Recall:} say the Hindi for an article, then the Hindi for an essay, then say \emph{chā̃ṭnā} again
\end{itemize}

\begin{tcolorbox}[breakable,colback=teal!4,colframe=teal!35!black,title={Before you move on}]
What does \hi{छाँटना} mean? (\textquotedblleft{}To sort out, to pick out\textquotedblright{}.) Say it once more.
\end{tcolorbox}
